% !TeX program = pdflatex
% Regulamin programu motywacyjnego (seria P) --- Basilisk Systems P.S.A.
% Regulamin, o którym mowa w § 7 ust. 2 i 5 umowy spółki (psa.tex). Odesłania pobierane z psa.aux i basilisk_sha.aux (xr-hyper).
\documentclass[11pt,a4paper]{report}
\usepackage[T1]{fontenc}
\usepackage[utf8]{inputenc}
\usepackage[english,polish]{babel}
\usepackage{xr-hyper}
\usepackage{basilisk_i18n}
\usepackage{lmodern}
\usepackage{microtype}
\usepackage{setspace}
\usepackage{parskip}
\usepackage{needspace}
\usepackage{tabularx}
\externaldocument{psa}[psa.pdf]
\externaldocument{basilisk_sha}[basilisk_sha.pdf]

\BusinessPlanCompany{Basilisk Systems}
\BusinessPlanWatermarkText{Basilisk Systems}
\BusinessPlanWatermarkColor{black!6}
\BusinessPlanWatermarkFontSize{38}{46}
\BusinessPlanWatermarkAngle{38}
\BusinessPlanAuthorsText{Basilisk Systems --- Walne Zgromadzenie}
\BusinessPlanConfidentialText{Projekt --- poufne}
\BusinessPlanFooterText{Basilisk Systems --- regulamin programu motywacyjnego (projekt)}
\BusinessPlanVersionText{Wersja 1.0 --- projekt}

\setstretch{1.1}
\setlist[enumerate]{topsep=0.2em,itemsep=0.15em,parsep=0.05em,leftmargin=2.0em}
\setlength{\parindent}{0pt}
\setlength{\parskip}{0.5em}
\setlength{\emergencystretch}{4em}
\renewcommand{\arraystretch}{1.2}

\newcounter{paragraf}
\newcommand{\paragraf}[1]{%
  \refstepcounter{paragraf}%
  \Needspace{5\baselineskip}%
  \subsection*{\S~\theparagraf\ #1}%
  \addcontentsline{toc}{section}{\S~\theparagraf\ #1}%
}
\newenvironment{ContractClause}[1]{\paragraf{#1}}{\par}
\newcommand{\field}[1]{\textbf{[#1]}}
\newcommand{\CompanyFull}{Basilisk Systems prosta spółka akcyjna}
\newcommand{\US}[1]{§~#1 umowy spółki}
\newcommand{\UA}[1]{§~#1 umowy akcjonariuszy}

\hypersetup{
  pdftitle={Basilisk Systems P.S.A. --- regulamin programu motywacyjnego (seria P)},
  pdfauthor={Basilisk Systems}
}

\begin{document}

\BusinessPlanTitlePage
  {Basilisk Systems}
  {Wersja 1.0 --- projekt}
  {Regulamin programu motywacyjnego}
  {Program motywacyjny Basilisk Systems P.S.A. --- akcje serii P}
  {Do umowy spółki Basilisk Systems P.S.A. w wersji 3.4-S}
  {1.0}

\tableofcontents
\newpage

\chapter*{Wprowadzenie}
\addcontentsline{toc}{chapter}{Wprowadzenie}

Regulamin jest regulaminem programu motywacyjnego, o którym mowa w \US{\ref{art:seria-p} ust.~2 i~5}. Przyjmuje go Walne Zgromadzenie uchwałą według projektu z Załącznika nr~3.

\textbf{Jak działa program:}
\begin{enumerate}
  \item Walne Zgromadzenie podejmuje uchwałę o warunkowej emisji akcji serii~P i wyraża zgodę na zawieranie umów uczestnictwa (art.~300$^{114}$ §~2 pkt~2 i §~3 KSH). Uchwała zmienia umowę spółki, więc wymaga protokołu notarialnego i wpisu do rejestru.
  \item Po wpisie Rada Dyrektorów zawiera z uczestnikami umowy uczestnictwa. Umowa daje prawo do objęcia określonej liczby akcji serii~P po cenie wykonania (\textbf{Opcje}).
  \item Opcje nabywa się stopniowo przez 4 lata. Po pierwszym roku nabywa się 25\% Opcji, a następnie co miesiąc 1/48.
  \item Nabyte Opcje uczestnik wykonuje w oknach wykonania, przez pisemne oświadczenie o objęciu akcji i zapłatę ceny wykonania. Akcje powstają z chwilą wpisu do rejestru akcjonariuszy.
\end{enumerate}

\textbf{Uwagi przed przyjęciem:}
\begin{enumerate}
  \item Uchwała o warunkowej emisji musi spełniać warunki art.~300$^{106}$ §~2 KSH (art.~300$^{115}$ §~2 KSH). Wymaga więc większości 4/5 głosów i uzasadnienia interesem Spółki, mimo że umowa spółki wyłącza prawo poboru akcji serii~P (\US{\ref{art:emisje} ust.~6}). Uchwała musi mieć także większość 75\% wszystkich głosów (\US{\ref{art:wz-zastrzezone} ust.~1}).
  \item Prawa do objęcia akcji przyznaje się dopiero po wpisie do rejestru zmiany umowy w przedmiocie warunkowej emisji (art.~300$^{116}$ §~2 KSH). Umowa uczestnictwa zawarta wcześniej jest tylko przyrzeczeniem.
  \item Odroczenie opodatkowania na podstawie art.~24 ust.~11--12b ustawy o PIT dotyczy programów utworzonych uchwałą walnego zgromadzenia spółki akcyjnej. Zastosowanie tych przepisów do prostej spółki akcyjnej trzeba potwierdzić interpretacją indywidualną przed pierwszym przydziałem \field{do weryfikacji przez doradcę podatkowego}.
\end{enumerate}

% ================================================================
\chapter*{Regulamin programu motywacyjnego}
\addcontentsline{toc}{chapter}{Regulamin programu motywacyjnego}

\begin{ContractClause}{Cel i definicje}
\label{esop:definicje}
\begin{enumerate}
  \item Celem programu jest pozyskanie i utrzymanie osób, od których zależy rozwój Spółki, oraz powiązanie ich wynagrodzenia ze wzrostem wartości Spółki.
  \item Użyte w regulaminie określenia oznaczają:
  \begin{enumerate}[label=\alph*)]
    \item \textbf{Pula} --- nie więcej niż 176 471 akcji serii~P (\US{\ref{art:emisje} ust.~3 lit.~b} i \US{\ref{art:seria-p} ust.~1});
    \item \textbf{Uczestnik} --- osoba, z którą Spółka zawarła Umowę Uczestnictwa;
    \item \textbf{Umowa Uczestnictwa} --- umowa przyznająca prawo do objęcia akcji serii~P (art.~300$^{114}$ §~2 pkt~2 KSH), według wzoru z Załącznika nr~1;
    \item \textbf{Opcja} --- prawo do objęcia jednej akcji serii~P po Cenie Wykonania;
    \item \textbf{Dzień Przyznania} --- dzień zawarcia Umowy Uczestnictwa albo wcześniejszy dzień wskazany w niej jako początek nabywania Opcji, nie wcześniejszy niż dzień rozpoczęcia współpracy;
    \item \textbf{Cena Wykonania} --- cena emisyjna akcji serii~P obejmowanej w wykonaniu Opcji, określona w Umowie Uczestnictwa;
    \item \textbf{Współpraca} --- stosunek pracy, umowa o świadczenie usług, pełnienie funkcji dyrektora albo inna stała współpraca Uczestnika ze Spółką lub jej spółką zależną;
    \item \textbf{Zdarzenie Wyjścia} --- nabycie przez osobę albo grupę osób działających w porozumieniu, niepowiązanych z Założycielami, akcji dających ponad 50\% głosów w Spółce, zbycie przedsiębiorstwa Spółki albo jego zasadniczej części, a także dopuszczenie akcji Spółki do obrotu na rynku regulowanym albo w alternatywnym systemie obrotu.
  \end{enumerate}
\end{enumerate}
\end{ContractClause}

\begin{ContractClause}{Podstawa i organy}
\begin{enumerate}
  \item Program prowadzi się na podstawie \US{\ref{art:seria-p}}, uchwały Walnego Zgromadzenia o warunkowej emisji akcji serii~P oraz niniejszego regulaminu, w trybie art.~300$^{114}$--300$^{118}$ KSH.
  \item Rada Dyrektorów prowadzi program: rekomenduje Uczestników, zawiera Umowy Uczestnictwa, prowadzi ewidencję Opcji, przyjmuje oświadczenia o objęciu akcji, wydaje dyspozycje wpisu do rejestru akcjonariuszy i dokonuje zgłoszeń rejestrowych (art.~300$^{117}$ §~2 i art.~300$^{118}$ §~2 KSH).
  \item Rada Dyrektorów przyznaje Opcje w granicach rocznego planu przydziałów zatwierdzonego przez Walne Zgromadzenie bezwzględną większością głosów oddanych. Plan określa liczbę Opcji dla kategorii osób i zasady ustalania Ceny Wykonania.
  \item Przyznanie Opcji dyrektorowi, akcjonariuszowi posiadającemu co najmniej 5\% akcji albo osobie im bliskiej wymaga uchwały Walnego Zgromadzenia. Umowę Uczestnictwa z dyrektorem zawiera w imieniu Spółki pełnomocnik powołany uchwałą akcjonariuszy (art.~300$^{79}$ §~1 KSH, \US{\ref{art:reprezentacja} ust.~3}).
  \item Liczba Opcji przyznanych i niewygasłych nie może przekroczyć Puli. Opcje wygasłe wracają do Puli i mogą zostać przyznane ponownie.
\end{enumerate}
\end{ContractClause}

\begin{ContractClause}{Uczestnicy}
\begin{enumerate}
  \item Uczestnikami mogą być pracownicy, współpracownicy, dyrektorzy i doradcy Spółki oraz jej spółek zależnych.
  \item Założyciele Pierwotni nie uczestniczą w programie, chyba że Walne Zgromadzenie postanowi inaczej większością 75\% wszystkich głosów.
  \item Przed zawarciem Umowy Uczestnictwa Rada Dyrektorów sprawdza, czy przyznanie Opcji i wydanie akcji nie wymaga zgody albo zawiadomienia na podstawie przepisów o kontroli inwestycji, kontroli eksportu i sankcjach oraz warunków koncesji Spółki (\US{\ref{art:zbywanie} ust.~2}).
\end{enumerate}
\end{ContractClause}

\begin{ContractClause}{Przyznanie Opcji}
\begin{enumerate}
  \item Opcje przyznaje się wyłącznie w Umowie Uczestnictwa zawartej w formie pisemnej albo elektronicznej (kwalifikowany podpis elektroniczny), po wpisie do rejestru zmiany umowy spółki w przedmiocie warunkowej emisji akcji serii~P.
  \item Umowa Uczestnictwa określa: liczbę Opcji, Dzień Przyznania, Cenę Wykonania, harmonogram nabywania, jeżeli odbiega od §~\ref{esop:nabywanie}, oraz termin wykonania.
  \item Cena Wykonania nie może być niższa niż 0,01 zł (\US{\ref{art:emisje} ust.~7}). Ustala się ją według zasad z planu przydziałów, w szczególności na poziomie wartości godziwej akcji w Dniu Przyznania albo na poziomie obniżonym w stosunku do ceny z ostatniej emisji akcji. Uzasadnienie zachowuje się w ewidencji.
  \item Opcje są nieodpłatne, niezbywalne i nie mogą być obciążone, z wyjątkiem dziedziczenia (§~\ref{esop:ustanie} ust.~5). Opcja nie daje praw akcjonariusza do chwili wpisu akcji do rejestru akcjonariuszy (art.~300$^{118}$ §~1 KSH).
\end{enumerate}
\end{ContractClause}

\begin{ContractClause}{Nabywanie Opcji}
\label{esop:nabywanie}
\begin{enumerate}
  \item Opcje nabywa się w okresie 48 miesięcy od Dnia Przyznania, pod warunkiem trwania Współpracy:
  \begin{enumerate}[label=\alph*)]
    \item po 12 miesiącach --- 25\% Opcji;
    \item po każdym kolejnym pełnym miesiącu --- 1/48 Opcji.
  \end{enumerate}
  Ułamki zaokrągla się w dół, a ostatnia transza wyrównuje różnicę.
  \item Nabywanie Opcji trwa w czasie choroby, urlopu macierzyńskiego, ojcowskiego, rodzicielskiego i wychowawczego oraz innych usprawiedliwionych nieobecności. Zawiesza się je na czas urlopu bezpłatnego albo przerwy we Współpracy uzgodnionej na wniosek Uczestnika, jeżeli trwają ponad 30 dni.
  \item Zmiana formy Współpracy albo przejście do spółki zależnej nie przerywa nabywania Opcji.
  \item Rada Dyrektorów może przyspieszyć nabywanie Opcji Uczestnika, z wyjątkiem Uczestnika będącego dyrektorem, któremu przyspieszenie przyznaje Walne Zgromadzenie.
\end{enumerate}
\end{ContractClause}

\begin{ContractClause}{Wykonanie Opcji}
\label{esop:wykonanie}
\begin{enumerate}
  \item Nabyte Opcje wykonuje się w oknach wykonania:
  \begin{enumerate}[label=\alph*)]
    \item w ciągu 30 dni od zatwierdzenia sprawozdania finansowego za każdy rok obrotowy;
    \item przed Zdarzeniem Wyjścia, w terminie wskazanym przez Radę Dyrektorów, nie krótszym niż 14 dni;
    \item w przypadkach, o których mowa w §~\ref{esop:ustanie}.
  \end{enumerate}
  \item Uczestnik wykonuje Opcje przez złożenie Spółce pisemnego oświadczenia o objęciu akcji według wzoru z Załącznika nr~2 (art.~300$^{117}$ §~1 KSH) i wpłatę Ceny Wykonania na rachunek Spółki. Oświadczenie złożone przed Zdarzeniem Wyjścia może być warunkowe.
  \item Niezwłocznie po otrzymaniu oświadczenia i wpłaty Rada Dyrektorów wydaje dyspozycję wpisu akcji do rejestru akcjonariuszy (art.~300$^{117}$ §~2 KSH). W terminie 30 dni po upływie roku kalendarzowego zgłasza do rejestru przedsiębiorców wykaz akcji objętych w tym roku (art.~300$^{118}$ §~2--3 KSH).
  \item Opcje wygasają z upływem 10 lat od Dnia Przyznania, a w każdym razie z dniem 30 listopada 2036 r., aby akcje zostały objęte i zgłoszone do rejestru przed upływem terminu z \US{\ref{art:emisje} ust.~3}.
  \item W razie Zdarzenia Wyjścia Rada Dyrektorów może, za zgodą Uczestnika, zastąpić wykonanie Opcji wypłatą pieniężną równą różnicy między ceną za akcję w transakcji a Ceną Wykonania, pomniejszoną o należne podatki i składki.
\end{enumerate}
\end{ContractClause}

\begin{ContractClause}{Ustanie Współpracy}
\label{esop:ustanie}
\begin{enumerate}
  \item Z dniem ustania Współpracy Opcje nienabyte wygasają.
  \item Opcje nabyte Uczestnik może wykonać w ciągu 6 miesięcy od ustania Współpracy; jeżeli w tym okresie nie przypada okno wykonania, Rada Dyrektorów wyznacza je w tym okresie.
  \item Wszystkie Opcje, także nabyte, wygasają, jeżeli Współpraca ustała z Ważnej Przyczyny. Ważną Przyczyną jest wyłącznie:
  \begin{enumerate}[label=\alph*)]
    \item rozwiązanie umowy z Uczestnikiem bez wypowiedzenia z jego winy (art.~52 §~1 Kodeksu pracy) albo z powodu rażącego naruszenia przez niego umowy cywilnoprawnej lub obowiązków dyrektora;
    \item prawomocne skazanie Uczestnika za umyślne przestępstwo na szkodę Spółki;
    \item istotne naruszenie przez Uczestnika obowiązku poufności, zakazu konkurencji albo obowiązków dotyczących własności intelektualnej, którego nie usunął w terminie 30 dni od wezwania, jeżeli usunięcie jest możliwe.
  \end{enumerate}
  Ciężar dowodu Ważnej Przyczyny spoczywa na Spółce.
  \item Jeżeli w ciągu 12 miesięcy po Zdarzeniu Wyjścia Współpraca ustanie z przyczyn innych niż Ważna Przyczyna albo rezygnacja Uczestnika, wszystkie Opcje Uczestnika uważa się za nabyte (podwójny warunek).
  \item W razie śmierci Uczestnika jego spadkobiercy mogą wykonać Opcje nabyte do dnia śmierci w ciągu 12 miesięcy od stwierdzenia nabycia spadku; wykonanie przez kilku spadkobierców wymaga wspólnego przedstawiciela.
\end{enumerate}
\end{ContractClause}

\begin{ContractClause}{Akcje serii P}
\begin{enumerate}
  \item Akcje serii~P są akcjami zwykłymi. Do ich zbywania stosuje się \US{\ref{art:zbywanie}}.
  \item Przed złożeniem oświadczenia o objęciu akcji Uczestnik przystępuje do umowy akcjonariuszy w zakresie określonym w \UA{\ref{sha:przystapienie} ust.~2}, w szczególności co do poufności, prawa pierwszeństwa, prawa przyłączenia i przymusowego współzbycia (\UA{\ref{sha:runda}}). Bez przystąpienia oświadczenie jest nieskuteczne.
  \item Akcje objęte w oknie, o którym mowa w §~\ref{esop:wykonanie} ust.~1 lit.~a, Uczestnik może zbyć nie wcześniej niż po 12 miesiącach od ich objęcia, chyba że Walne Zgromadzenie postanowi inaczej albo nastąpi Zdarzenie Wyjścia.
\end{enumerate}
\end{ContractClause}

\begin{ContractClause}{Zmiany w kapitale i korekty}
\begin{enumerate}
  \item W razie podziału albo scalenia akcji, emisji akcji dla dotychczasowych akcjonariuszy bez wkładów albo innej zmiany o podobnym skutku Rada Dyrektorów proporcjonalnie koryguje liczbę Opcji i Cenę Wykonania, tak aby wartość Opcji nie uległa zmianie. Korekta wymagająca zwiększenia Puli wymaga zmiany umowy spółki (\US{\ref{art:seria-p} ust.~3}).
  \item W razie połączenia, podziału albo przekształcenia Spółki Opcje przechodzą na prawa równoważne w spółce przejmującej albo przekształconej, jeżeli plan połączenia, podziału albo przekształcenia tak stanowi; w przeciwnym razie stosuje się §~\ref{esop:wykonanie} ust.~1 lit.~b i ust.~5.
\end{enumerate}
\end{ContractClause}

\begin{ContractClause}{Podatki}
\begin{enumerate}
  \item Uczestnik samodzielnie rozlicza podatki związane z Opcjami i akcjami, z zastrzeżeniem obowiązków Spółki jako płatnika.
  \item Spółka dąży do tego, aby program spełniał warunki art.~24 ust.~11--12b ustawy o PIT, ale nie gwarantuje określonych skutków podatkowych. Przed pierwszym przydziałem Spółka występuje o interpretację indywidualną i udostępnia ją Uczestnikom \field{do weryfikacji przez doradcę podatkowego}.
\end{enumerate}
\end{ContractClause}

\begin{ContractClause}{Postanowienia końcowe}
\begin{enumerate}
  \item Zmiana regulaminu wymaga uchwały Walnego Zgromadzenia podjętej większością 75\% wszystkich głosów. Zmiana pogarszająca warunki Opcji już przyznanych wymaga zgody Uczestnika.
  \item Spółka przetwarza dane osobowe Uczestników w celu zawarcia i wykonania Umowy Uczestnictwa (art.~6 ust.~1 lit.~b RODO) oraz wykonania obowiązków prawnych (art.~6 ust.~1 lit.~c RODO).
  \item Informacje o programie, Cenie Wykonania i wartości Spółki są informacjami poufnymi w rozumieniu polityki bezpieczeństwa informacji.
  \item Uczestnictwo w programie nie stanowi składnika wynagrodzenia za pracę i nie daje roszczenia o kolejne przydziały.
  \item Do regulaminu i Umów Uczestnictwa stosuje się prawo polskie. Spory rozstrzyga sąd powszechny właściwy dla siedziby Spółki; spory ze stosunku pracy --- sąd pracy.
  \item Regulamin sporządzono w języku polskim i angielskim; w razie rozbieżności rozstrzyga wersja polska. Regulamin wchodzi w życie z dniem wpisu do rejestru zmiany umowy spółki w przedmiocie warunkowej emisji akcji serii~P.
\end{enumerate}
\end{ContractClause}

% ================================================================
\chapter*{Załącznik nr 1 --- Wzór Umowy Uczestnictwa}
\addcontentsline{toc}{chapter}{Załącznik nr 1 --- Wzór Umowy Uczestnictwa}

zawarta w dniu \field{data} pomiędzy \textbf{\CompanyFull{}} z siedzibą w Pruszkowie, KRS \field{numer}, reprezentowaną przez \field{dwóch dyrektorów albo pełnomocnika z uchwały akcjonariuszy --- jeżeli Uczestnik jest dyrektorem} (\textbf{„Spółka”}), a \field{imię, nazwisko, PESEL, adres} (\textbf{„Uczestnik”}).

\begin{enumerate}
  \item Spółka przyznaje Uczestnikowi \field{liczba} Opcji, z których każda uprawnia do objęcia jednej akcji serii~P po Cenie Wykonania \field{kwota} zł, na zasadach regulaminu programu motywacyjnego przyjętego uchwałą Walnego Zgromadzenia nr \field{numer} z dnia \field{data} (\textbf{„Regulamin”}), który stanowi część umowy.
  \item Dzień Przyznania: \field{data}. Harmonogram nabywania: zgodnie z §~\ref{esop:nabywanie} Regulaminu \field{albo: odmienny harmonogram}.
  \item Prawo do objęcia akcji przyznaje się na podstawie art.~300$^{114}$ §~2 pkt~2 KSH, za zgodą Walnego Zgromadzenia wyrażoną w uchwale nr \field{numer} z dnia \field{data}. Zmianę umowy spółki w przedmiocie warunkowej emisji wpisano do rejestru w dniu \field{data}.
  \item Uczestnik oświadcza, że zapoznał się z Regulaminem, umową spółki, umową akcjonariuszy i polityką bezpieczeństwa informacji oraz że nie jest objęty sankcjami.
  \item Umowa nie zmienia warunków Współpracy Uczestnika i nie stanowi gwarancji jej trwania.
  \item Zmiany umowy wymagają formy pisemnej albo elektronicznej pod rygorem nieważności.
\end{enumerate}

\vspace{1em}
\field{za Spółkę} \hfill \field{Uczestnik}

% ================================================================
\chapter*{Załącznik nr 2 --- Wzór oświadczenia o objęciu akcji serii P}
\addcontentsline{toc}{chapter}{Załącznik nr 2 --- Wzór oświadczenia o objęciu akcji serii P}

Ja, \field{imię, nazwisko, PESEL, adres, adres e-mail do rejestru akcjonariuszy}, działając na podstawie art.~300$^{117}$ §~1 KSH, uchwały Walnego Zgromadzenia \CompanyFull{} nr \field{numer} z dnia \field{data} o warunkowej emisji akcji serii~P oraz Umowy Uczestnictwa z dnia \field{data}, oświadczam, że:
\begin{enumerate}
  \item obejmuję \field{liczba} akcji zwykłych serii~P po cenie emisyjnej \field{kwota} zł za akcję, łącznie \field{kwota} zł, którą wpłaciłem/wpłaciłam w dniu \field{data} na rachunek Spółki;
  \item \field{jeżeli dotyczy:} oświadczenie składam pod warunkiem zawarcia transakcji stanowiącej Zdarzenie Wyjścia do dnia \field{data};
  \item przystąpiłem/przystąpiłam do umowy akcjonariuszy w zakresie wymaganym przez Regulamin;
  \item proszę o wpis akcji do rejestru akcjonariuszy na moją rzecz.
\end{enumerate}

\field{miejscowość, data} \hfill \field{podpis własnoręczny albo kwalifikowany podpis elektroniczny}

% ================================================================
\chapter*{Załącznik nr 3 --- Projekt uchwały Walnego Zgromadzenia}
\addcontentsline{toc}{chapter}{Załącznik nr 3 --- Projekt uchwały Walnego Zgromadzenia}

\textbf{Uchwała nr \field{numer} Walnego Zgromadzenia \CompanyFull{} z dnia \field{data} w sprawie warunkowej emisji akcji serii~P, zmiany umowy spółki, przyjęcia regulaminu programu motywacyjnego i zgody na zawieranie umów uczestnictwa}

{\small (protokół notarialny; większość co najmniej 4/5 głosów, art.~300$^{115}$ §~2 w zw. z art.~300$^{106}$ §~2 KSH, i nie mniej niż 75\% wszystkich głosów, \US{\ref{art:wz-zastrzezone} ust.~1})\par}

\begin{enumerate}
  \item Na podstawie art.~300$^{114}$ §~1 i §~2 pkt~2 KSH oraz \US{\ref{art:seria-p}} Walne Zgromadzenie postanawia o warunkowej emisji nie więcej niż \field{liczba, nie więcej niż 176 471} akcji zwykłych serii~P.
  \item Cena emisyjna akcji serii~P jest równa Cenie Wykonania określonej w umowie uczestnictwa według zasad regulaminu programu motywacyjnego, nie niższa niż 0,01 zł.
  \item Celem warunkowej emisji jest przyznanie praw do objęcia akcji serii~P uczestnikom programu motywacyjnego Spółki.
  \item Prawo do objęcia akcji serii~P można wykonać do dnia 30 listopada 2036 r.
  \item Osobami uprawnionymi do objęcia akcji serii~P są pracownicy, współpracownicy, dyrektorzy i doradcy Spółki oraz jej spółek zależnych, z którymi Spółka zawrze umowy uczestnictwa.
  \item W interesie Spółki --- w celu pozyskania i utrzymania kluczowych osób --- akcje serii~P obejmowane są z wyłączeniem prawa poboru; umowa spółki wyłącza to prawo (\US{\ref{art:emisje} ust.~6}). \field{uzasadnienie i sposób ustalenia ceny emisyjnej}.
  \item Walne Zgromadzenie zmienia umowę spółki przez dodanie w \US{\ref{art:seria-p}} ustępu w brzmieniu: „Na podstawie uchwały Walnego Zgromadzenia nr \field{numer} z dnia \field{data} dokonuje się warunkowej emisji nie więcej niż \field{liczba} akcji zwykłych serii~P w celu przyznania praw do ich objęcia uczestnikom programu motywacyjnego; prawo do objęcia akcji można wykonać do dnia 30 listopada 2036 r.”
  \item Na podstawie art.~300$^{114}$ §~3 KSH Walne Zgromadzenie wyraża zgodę na zawieranie przez Spółkę umów uczestnictwa według wzoru stanowiącego Załącznik nr~1 do regulaminu, w granicach planu przydziałów.
  \item Walne Zgromadzenie przyjmuje regulamin programu motywacyjnego w brzmieniu stanowiącym załącznik do uchwały oraz plan przydziałów na rok \field{rok}.
  \item Uchwała wchodzi w życie z chwilą podjęcia, z tym że zmiana umowy spółki wywołuje skutek z chwilą wpisu do rejestru.
\end{enumerate}

\end{document}
